%%%%%%%%%%%%%%%%%%%%%%%%%%%%%%%%%%%%%%%%%%%%%%%%%%%%%%%%%%%
%% Begin exercise %%
%%%%%%%%%%%%%%%%%%%%%%%%%%%%%%%%%%%%%%%%%%%%%%%%%%%%%%%%%%%
\ex{Capacitors in Power Electronics}

%%%%%%%%%%%%%%%%%%%%%%%%%%%%%%%%%%%%%%%%%%%%%%%%%%%%%%%%%%%%%%%%%%%%%%%%%%%%%%%
%% Task 1 Why one capacitor cannot perform every function
%%%%%%%%%%%%%%%%%%%%%%%%%%%%%%%%%%%%%%%%%%%%%%%%%%%%%%%%%%%%%%%%%%%%%%%%%%%%%%%

\task{Why one capacitor cannot perform every function}

A \SI{400}{\volt} DC bus uses the three capacitor branches listed in
\autoref{table:capacitor_task1_data}. The MLCC bank has a nominal capacitance of
\SI{4.7}{\micro\farad}, but its effective capacitance falls to
\SI{2.2}{\micro\farad} under the applied DC bias.

\begin{table}[ht]
    \centering
    \begin{tabular}{llll}
        \toprule
        Technology & Effective capacitance & ESR & ESL \\
        \midrule
        Aluminum electrolytic & \SI{1000}{\micro\farad} & \SI{40}{\milli\ohm} & \SI{20}{\nano\henry} \\
        Film & \SI{20}{\micro\farad} & \SI{8}{\milli\ohm} & \SI{8}{\nano\henry} \\
        MLCC bank & \SI{2.2}{\micro\farad} & \SI{3}{\milli\ohm} & \SI{1}{\nano\henry} \\
        \bottomrule
    \end{tabular}
    \caption{Capacitor data for Task 1.}
    \label{table:capacitor_task1_data}
\end{table}

Use
\begin{align*}
    E_C &= \frac{1}{2}CV^2, \quad 
    f_\mathrm{SRF} = \frac{1}{2\pi\sqrt{L_\mathrm{ESL}C}} \quad  \text{and} \\
    Z(j\omega) &= R_\mathrm{ESR} + j\left(\omega L_\mathrm{ESL}-\frac{1}{\omega C}\right).
\end{align*}

\subtask{Calculate the energy stored in each capacitor at \SI{400}{\volt}.}

\begin{solutionblock}
    For the electrolytic capacitor,
    \begin{equation}
        E_\mathrm{el}=\frac{1}{2}\cdot\SI{1000}{\micro\farad}\cdot(\SI{400}{\volt})^2
        =\SI{80}{\joule}.
    \end{equation}
    For the film capacitor,
    \begin{equation}
        E_\mathrm{film}=\frac{1}{2}\cdot\SI{20}{\micro\farad}\cdot(\SI{400}{\volt})^2
        =\SI{1.6}{\joule}.
    \end{equation}
    For the MLCC bank,
    \begin{equation}
        E_\mathrm{MLCC}=\frac{1}{2}\cdot\SI{2.2}{\micro\farad}\cdot(\SI{400}{\volt})^2
        =\SI{0.176}{\joule}.
    \end{equation}
    The electrolytic capacitor dominates the energy-buffering function.
\end{solutionblock}

\subtask{Calculate the self-resonant frequency of each capacitor.}

\begin{solutionblock}
    \begin{align}
        f_\mathrm{SRF,el}
        &=\frac{1}{2\pi\sqrt{\SI{20}{\nano\henry}\cdot\SI{1000}{\micro\farad}}}
        \approx\SI{35.6}{\kilo\hertz},\\
        f_\mathrm{SRF,film}
        &=\frac{1}{2\pi\sqrt{\SI{8}{\nano\henry}\cdot\SI{20}{\micro\farad}}}
        \approx\SI{398}{\kilo\hertz},\\
        f_\mathrm{SRF,MLCC}
        &=\frac{1}{2\pi\sqrt{\SI{1}{\nano\henry}\cdot\SI{2.2}{\micro\farad}}}
        \approx\SI{3.39}{\mega\hertz}.
    \end{align}
\end{solutionblock}

\subtask{Calculate the approximate impedance magnitude of each capacitor at
\SI{100}{\hertz}, \SI{300}{\kilo\hertz}, and \SI{5}{\mega\hertz}. Identify the
most effective branch at each frequency.}

\begin{solutionblock}
    The impedance magnitude is
    \begin{equation}
        |Z|=\sqrt{R_\mathrm{ESR}^2+
        \left(2\pi fL_\mathrm{ESL}-\frac{1}{2\pi fC}\right)^2}.
    \end{equation}

    \begin{center}
        \begin{tabular}{llll}
            \toprule
            Frequency & Electrolytic & Film & MLCC \\
            \midrule
            \SI{100}{\hertz} & \SI{1.59}{\ohm} & \SI{79.6}{\ohm} & \SI{723}{\ohm} \\
            \SI{300}{\kilo\hertz} & \SI{54.6}{\milli\ohm} & \SI{14.0}{\milli\ohm} & \SI{239}{\milli\ohm} \\
            \SI{5}{\mega\hertz} & \SI{630}{\milli\ohm} & \SI{250}{\milli\ohm} & \SI{17.2}{\milli\ohm} \\
            \bottomrule
        \end{tabular}
    \end{center}

    Therefore, the electrolytic branch is most effective at \SI{100}{\hertz},
    the film branch at \SI{300}{\kilo\hertz}, and the MLCC branch at
    \SI{5}{\mega\hertz}. Nominal capacitance alone is therefore not a sufficient
    selection criterion.
\end{solutionblock}

\subtask{The switching current changes by \SI{30}{\ampere} within
\SI{10}{\nano\second}. The total loop inductance is \SI{12}{\nano\henry}.
Estimate the resulting voltage spike.}

\begin{solutionblock}
    \begin{align}
        \frac{\mathrm{d}i}{\mathrm{d}t}
        &=\frac{\SI{30}{\ampere}}{\SI{10}{\nano\second}}
        =\SI{3}{\ampere\per\nano\second},\\
        v_\mathrm{spike}
        &=L_\mathrm{loop}\frac{\mathrm{d}i}{\mathrm{d}t}
        =\SI{12}{\nano\henry}\cdot\SI{3}{\ampere\per\nano\second}
        =\SI{36}{\volt}.
    \end{align}
    A low-ESL capacitor cannot compensate for a poorly designed external
    commutation loop.
\end{solutionblock}

\subtask{The electrolytic capacitor has an insulation resistance of
\SI{20}{\mega\ohm}. Calculate its leakage current and leakage power at
\SI{400}{\volt}.}

\begin{solutionblock}
    \begin{align}
        I_\mathrm{leak}&=\frac{\SI{400}{\volt}}{\SI{20}{\mega\ohm}}
        =\SI{20}{\micro\ampere},\\
        P_\mathrm{leak}&=VI=\SI{400}{\volt}\cdot\SI{20}{\micro\ampere}
        =\SI{8}{\milli\watt}.
    \end{align}
\end{solutionblock}

\subtask{Explain why all three capacitor technologies may be required on the same DC bus.}

\begin{solutionblock}
    The electrolytic branch supplies low-frequency energy, the film branch
    carries substantial mid-frequency ripple current with low loss, and the
    MLCC branch provides the local high-frequency switching-current path.
    Their parallel combination must nevertheless be checked for anti-resonance
    and layout-induced impedance peaks.
\end{solutionblock}


%%%%%%%%%%%%%%%%%%%%%%%%%%%%%%%%%%%%%%%%%%%%%%%%%%%%%%%%%%%%%%%%%%%%%%%%%%%%%%%
%% Task 2 Input and output capacitor co-design for a buck converter
%%%%%%%%%%%%%%%%%%%%%%%%%%%%%%%%%%%%%%%%%%%%%%%%%%%%%%%%%%%%%%%%%%%%%%%%%%%%%%%


\task{Input and output capacitor co-design for a buck converter}

A buck converter operates with
\begin{equation*}
    V_\mathrm{in}=\SI{48}{\volt},\qquad
    V_o=\SI{12}{\volt},\qquad
    I_o=\SI{10}{\ampere},\qquad
    f_s=\SI{200}{\kilo\hertz}.
\end{equation*}
The allowable peak-to-peak input ripple is $\SI{0.20}{\volt}$. The output
inductor ripple is $\SI{4}{\ampere}$. A candidate output capacitor has an
effective capacitance of $\SI{100}{\micro\farad}$ and an ESR of
$\SI{5}{\milli\ohm}$. \newline

For a load transient, the load current increases suddenly by \SI{8}{\ampere}.
The output inductor is \SI{10}{\micro\henry}, and the maximum available voltage
for increasing its current is \SI{6}{\volt}. The maximum permitted output
voltage drop is \SI{0.20}{\volt}. \newline
Use $\quad C_\mathrm{in,min} \approx \frac{I_oD(1-D)}{f_s\Delta V_\mathrm{in,pp}}, \quad  
    \Delta v_\mathrm{C} \approx\frac{\Delta i_\mathrm{L}}{8f_sC_o},\quad \Delta v_\mathrm{ESR} = \Delta i_LR_\mathrm{ESR},  
    \quad \text{and} \quad Z_\mathrm{target}=\frac{\Delta V_\mathrm{max}}{\Delta I_\mathrm{load}}.$


\subtask{Calculate the converter duty ratio and the minimum effective input capacitance.}

\begin{solutionblock}
    \begin{align}
        D&=\frac{V_o}{V_\mathrm{in}}=\frac{\SI{12}{\volt}}{\SI{48}{\volt}}=0.25,\\
        C_\mathrm{in,min}
        &\approx \frac{I_oD(1-D)}{f_s\Delta V_\mathrm{in,pp}} = \frac{\SI{10}{\ampere}(0.25)(0.75)}
        {\SI{200}{\kilo\hertz} \cdot \SI{0.20}{\volt}}
        \approx\SI{46.9}{\micro\farad}.
    \end{align}
\end{solutionblock}

\subtask{Compare one nominal \SI{100}{\micro\farad} MLCC that retains
\SI{45}{\percent} of its capacitance with two nominal
\SI{47}{\micro\farad} MLCCs that each retain \SI{55}{\percent}.}

\begin{solutionblock}
    \begin{align}
        C_\mathrm{eff,100}&=0.45\cdot\SI{100}{\micro\farad}
        =\SI{45}{\micro\farad},\\
        C_\mathrm{eff,2x47}&=2\cdot0.55\cdot\SI{47}{\micro\farad}
        =\SI{51.7}{\micro\farad}.
    \end{align}
    The first choice is slightly below the required effective capacitance,
    whereas the two-capacitor solution meets the first-order requirement.
    Tolerance, ageing, and temperature derating still require additional margin.
\end{solutionblock}

\subtask{Calculate the ideal capacitive output ripple and the ESR ripple.}

\begin{solutionblock}
    \begin{align}
        \Delta v_\mathrm{C}&\approx \frac{\Delta i_\mathrm{L}}{8f_sC_o} = \frac{\SI{4}{\ampere}}
        {8\cdot\SI{200}{\kilo\hertz}\cdot\SI{100}{\micro\farad}}
        =\SI{25}{\milli\volt},\\
        \Delta v_\mathrm{ESR}&=\Delta i_\mathrm{L} R_\mathrm{ESR}=\SI{4}{\ampere}\cdot\SI{5}{\milli\ohm}
        =\SI{20}{\milli\volt}.
    \end{align}
    Neglecting ESL, the approximate total ripple is \SI{45}{\milli\volt}.
\end{solutionblock}

\subtask{During the \SI{8}{\ampere} load step, calculate the inductor current
slew rate, the response time, and the charge supplied by the output capacitor.}

\begin{solutionblock}
    \begin{align}
        \frac{\mathrm{d}i_\mathrm{L}}{\mathrm{d}t}
        &=\frac{V_\mathrm{L, max}}{L}=\frac{\SI{6}{\volt}}{\SI{10}{\micro\henry}}
        =\SI{0.6}{\ampere\per\micro\second},\\
        t_r&=\frac{\Delta I_\mathrm{load}}{\mathrm{d}i_\mathrm{L} / \mathrm{d}t}
            =\frac{\SI{8}{\ampere}}{\SI{0.6}{\ampere\per\micro\second}}
        \approx\SI{13.3}{\micro\second}.
    \end{align}
    The current deficit decreases approximately linearly, so
    \begin{equation}
        \Delta Q=\frac{1}{2}\Delta I_\mathrm{load}t_r
        =\frac{1}{2}\cdot\SI{8}{\ampere}\cdot\SI{13.3}{\micro\second}
        \approx\SI{53.3}{\micro\coulomb}.
    \end{equation}
\end{solutionblock}

\subtask{Determine the capacitance required to keep the total transient drop below
\SI{0.20}{\volt}.}

\begin{solutionblock}
    The immediate ESR drop is
    \begin{equation}
        \Delta v_\mathrm{ESR}=\Delta i_\mathrm{L} R_\mathrm{ESR}=\SI{8}{\ampere}\cdot\SI{5}{\milli\ohm}
        =\SI{40}{\milli\volt}.
    \end{equation}
    The remaining capacitive voltage budget is therefore
    \SI{0.16}{\volt}. Hence,
    \begin{equation}
        C_\mathrm{min}=\frac{\Delta Q}{\Delta V_\mathrm{C}}
        =\frac{\SI{53.3}{\micro\coulomb}}{\SI{0.16}{\volt}}
        \approx\SI{333}{\micro\farad}.
    \end{equation}
    The \SI{100}{\micro\farad} capacitor satisfies the steady-state ripple
    requirement but not the load-transient requirement.
\end{solutionblock}

\subtask{Calculate the target impedance and explain why ripple-based capacitance alone is insufficient.}

\begin{solutionblock}
    \begin{equation}
        Z_\mathrm{target}=\frac{\Delta V_\mathrm{max}}{\Delta I_\mathrm{load}}=\frac{\SI{0.20}{\volt}}{\SI{8}{\ampere}}
        =\SI{25}{\milli\ohm}.
    \end{equation}
    The output network must remain below this impedance over the frequency range
    where the control loop and inductor cannot yet restore the load current.
    This normally requires a coordinated bulk, polymer, film, and/or ceramic network.
\end{solutionblock}


%%%%%%%%%%%%%%%%%%%%%%%%%%%%%%%%%%%%%%%%%%%%%%%%%%%%%%%%%%%%%%%%%%%%%%%%%%%%%%%
%% Task 3 DC-link capacitor bank, current capability, and layout
%%%%%%%%%%%%%%%%%%%%%%%%%%%%%%%%%%%%%%%%%%%%%%%%%%%%%%%%%%%%%%%%%%%%%%%%%%%%%%%

\task{DC-link capacitor bank, current capability, and layout}

A converter must supply a temporary power mismatch of \SI{25}{\kilo\watt} for
\SI{5}{\milli\second}. The DC-link voltage may fall from \SI{800}{\volt} to
\SI{720}{\volt}. Available electrolytic capacitors have
\begin{equation*}
    C=\SI{4700}{\micro\farad},\quad
    V_\mathrm{rated}=\SI{450}{\volt},\quad
    R_\mathrm{ESR}=\SI{18}{\milli\ohm},\quad
    I_\mathrm{ripple,rated}=\SI{18}{\ampere}.
\end{equation*}
The required DC-link ripple-current capability is \SI{30}{\ampere}. Use
$\quad C_\mathrm{req}=\frac{2P \cdot t_\mathrm{h}}{V_\mathrm{max}^2-V_\mathrm{min}^2}.$


\subtask{Calculate the minimum capacitance from the hold-up requirement.}

\begin{solutionblock}
    \begin{equation}
        C_\mathrm{req} = C_\mathrm{req}=\frac{2P \cdot t_\mathrm{h}}{V_\mathrm{max}^2-V_\mathrm{min}^2}
        =\frac{2\cdot\SI{25}{\kilo\watt}\cdot\SI{5}{\milli\second}}
        {(\SI{800}{\volt})^2-(\SI{720}{\volt})^2}
        \approx\SI{2.06}{\milli\farad}.
    \end{equation}
\end{solutionblock}

\subtask{Evaluate one string of two series-connected capacitors in terms of
equivalent capacitance, voltage rating, ESR, and ripple-current capability.}

\begin{solutionblock}
    \begin{align}
        C_\mathrm{string}&=\frac{C}{2}=\frac{\SI{4700}{\micro\farad}}{2}=\SI{2350}{\micro\farad},\\
        V_\mathrm{string}&=2 \cdot V_\mathrm{rated} = 2\cdot\SI{450}{\volt}=\SI{900}{\volt},\\
        R_\mathrm{ESR,string}&=2 \cdot R_\mathrm{ESR} =2\cdot\SI{18}{\milli\ohm}=\SI{36}{\milli\ohm}.
    \end{align}
    The ripple-current rating remains approximately \SI{18}{\ampere} because
    the same current passes through both series capacitors. The string meets the
    capacitance and voltage requirements, but not the required ripple-current capability.
\end{solutionblock}

\subtask{Evaluate two identical series strings connected in parallel and determine
whether the resulting bank satisfies all first-order requirements.}

\begin{solutionblock}
    \begin{align}
        C_\mathrm{bank}&=2 \cdot C_\mathrm{string}=2\cdot\SI{2350}{\micro\farad}
        =\SI{4700}{\micro\farad},\\
        R_\mathrm{ESR,bank}&=\frac{R_\mathrm{ESR,string}}{2}=\frac{\SI{36}{\milli\ohm}}{2}
        =\SI{18}{\milli\ohm},\\
        I_\mathrm{ripple,bank}&\approx 2 \cdot I_\mathrm{ripple,rated} = 2 \cdot \SI{18}{\ampere}
        =\SI{36}{\ampere}.
    \end{align}
    The two-string bank satisfies the capacitance, voltage, and ripple-current requirements.
\end{solutionblock}

\subtask{Calculate the ESR loss of the complete bank at \SI{30}{\ampere} RMS,
assuming equal current sharing.}

\begin{solutionblock}
    \begin{equation}
        P_\mathrm{ESR}=I_\mathrm{rms}^2 R_\mathrm{ESR,bank}
        =(\SI{30}{\ampere})^2\cdot\SI{18}{\milli\ohm}
        =\SI{16.2}{\watt}.
    \end{equation}
    Each string carries \SI{15}{\ampere}; therefore, each of the four capacitors
    dissipates approximately
    \begin{equation}
        P_\mathrm{cap}= (\frac{I_\mathrm{rms}}{2})^2 R_\mathrm{ESR,bank}=(\SI{15}{\ampere})^2\cdot\SI{18}{\milli\ohm}
        =\SI{4.05}{\watt}.
    \end{equation}
\end{solutionblock}

\subtask{Explain why voltage-balancing resistors are needed.}

\begin{solutionblock}
    Capacitance tolerance, leakage-current mismatch, temperature, and ageing can
    produce unequal voltage sharing. One capacitor may therefore exceed its
    individual voltage rating even when the total string voltage is acceptable.
    Passive or active balancing is needed to control this imbalance.
\end{solutionblock}

\subtask{A local film/ceramic branch and laminated busbar reduce the total loop
inductance from \SI{15}{\nano\henry} to \SI{4}{\nano\henry}. Calculate the
switching spike before and after the improvement for
$\mathrm{d}i/\mathrm{d}t=\SI{3}{\ampere\per\nano\second}$.}

\begin{solutionblock}
    \begin{align}
        v_{\sigma,1}&=L_\mathrm{loop, org} \cdot \frac{\mathrm{d}i}{\mathrm{d}t}=\SI{15}{\nano\henry}\cdot
        \SI{3}{\ampere\per\nano\second}=\SI{45}{\volt},\\
        v_{\sigma,2}&=L_\mathrm{loop, improved} \cdot \frac{\mathrm{d}i}{\mathrm{d}t}=\SI{4}{\nano\henry}\cdot
        \SI{3}{\ampere\per\nano\second}=\SI{12}{\volt}.
    \end{align}
    The layout and local capacitor placement reduce the overshoot by \SI{33}{\volt}.
\end{solutionblock}

\subtask{Explain why a hybrid bank may improve electrolytic-capacitor lifetime.}

\begin{solutionblock}
    Film capacitors can carry much of the mid-frequency ripple current, while
    local ceramics carry the highest-frequency switching current. This reduces
    high-frequency heating in the electrolytic capacitors and can extend their
    useful life. Anti-resonance and unequal current sharing must still be checked.
\end{solutionblock}


%%%%%%%%%%%%%%%%%%%%%%%%%%%%%%%%%%%%%%%%%%%%%%%%%%%%%%%%%%%%%%%%%%%%%%%%%%%%%%%
%% Task 4 Why snubber, resonant, and EMI capacitors are not interchangeable
%%%%%%%%%%%%%%%%%%%%%%%%%%%%%%%%%%%%%%%%%%%%%%%%%%%%%%%%%%%%%%%%%%%%%%%%%%%%%%%

\task{Why snubber, resonant, and EMI capacitors are not interchangeable}

Three capacitor applications are being designed. \newline

\textbf{Case A -- WBG snubber:} \newline
$C_s=\SI{4.7}{\nano\farad}$ is connected across an \SI{800}{\volt} switch,
$\mathrm{d}v/\mathrm{d}t=\SI{20}{\volt\per\nano\second}$, and
$f_s=\SI{50}{\kilo\hertz}$. \newline

\textbf{Case B -- resonant capacitor:} \newline
$L_r=\SI{10}{\micro\henry}$, required $f_r=\SI{250}{\kilo\hertz}$, and
capacitance tolerance $\pm\SI{5}{\percent}$. \newline

\textbf{Case C -- Y capacitors:} \newline
A $\SI{230}{\volt}$ RMS, $\SI{50}{\hertz}$ supply permits a maximum total
line-to-earth leakage current of $\SI{0.5}{\milli\ampere}$ through two identical
Y capacitors.

\subtask{For the snubber, calculate the peak capacitor current, stored energy,
and an upper-bound repetitive power if the stored energy is dissipated once per cycle.}

\begin{solutionblock}
    \begin{align}
        I_\mathrm{pk}&=C_s\frac{\mathrm{d}v}{\mathrm{d}t}
        =\SI{4.7}{\nano\farad}\cdot\SI{20}{\volt\per\nano\second}
        =\SI{94}{\ampere},\\
        E_s&=\frac{1}{2}C_sV^2
        =\frac{1}{2}\cdot\SI{4.7}{\nano\farad}\cdot(\SI{800}{\volt})^2
        \approx\SI{1.50}{\milli\joule},\\
        P_\mathrm{upper}&=E_sf_s
        =\SI{1.50}{\milli\joule}\cdot\SI{50}{\kilo\hertz}
        \approx\SI{75.2}{\watt}.
    \end{align}
    The actual energy exchange depends on the snubber topology, but pulse-current,
    pulse-energy, and thermal ratings cannot be ignored.
\end{solutionblock}

\subtask{Calculate the resonant capacitance required in Case B and the frequency
range caused by the \SI{5}{\percent} capacitance tolerance.}

\begin{solutionblock}
    \begin{align}
        C_r&=\frac{1}{(2\pi f_r)^2L_r}
        \approx\SI{40.5}{\nano\farad},\\
        f_{r,\min}&=\frac{f_{r,\min}}{\sqrt{1+\Delta C/C}} = \frac{\SI{250}{\kilo\hertz}}{\sqrt{1.05}}
        \approx\SI{244.0}{\kilo\hertz},\\
        f_{r,\max}&=\frac{f_{r,\min}}{\sqrt{1-\Delta C/C}} = \frac{\SI{250}{\kilo\hertz}}{\sqrt{0.95}}
        \approx\SI{256.5}{\kilo\hertz}.
    \end{align}
    Thus, a \SI{5}{\percent} capacitance tolerance produces approximately a
    \SI{2.5}{\percent} resonant-frequency variation.
\end{solutionblock}

\subtask{Determine the maximum value of each Y capacitor in Case C. Check whether
two \SI{4.7}{\nano\farad} Y capacitors satisfy the leakage-current limit.}

\begin{solutionblock}
    The total leakage current through two equal capacitors is
    \begin{equation}
        I_\mathrm{leak,total}=2\omega C \cdot V.
    \end{equation}
    Hence,
    \begin{equation}
        C_\mathrm{max}=\frac{I_\mathrm{leak,total}}{2\omega \cdot  V}=\frac{\SI{0.5}{\milli\ampere}}
        {2\cdot2\pi \cdot \SI{50}{\hertz}\cdot\SI{230}{\volt}}
        \approx\SI{3.46}{\nano\farad}.
    \end{equation}
    A standard \SI{3.3}{\nano\farad} value satisfies the first-order limit.
    For two \SI{4.7}{\nano\farad} capacitors,
    \begin{equation}
        I_\mathrm{leak}=2\omega C \cdot V=2 \cdot2 \pi \cdot \SI{50}{\hertz} \cdot
        \SI{4.7}{\nano\farad} \cdot \SI{230}{\volt}
        \approx\SI{0.679}{\milli\ampere},
    \end{equation}
    which exceeds the permitted value.
\end{solutionblock}

\subtask{Recommend an appropriate capacitor technology for each case and explain
why nominal capacitance and voltage rating are not sufficient selection criteria.}

\begin{solutionblock}
    A pulse-rated low-ESL film or power-ceramic capacitor is appropriate for the
    snubber. A low-loss and capacitance-stable polypropylene film or suitable
    resonant ceramic capacitor is appropriate for the resonant tank. Only a
    certified Y-class safety capacitor is acceptable from line to earth.
    Pulse-current capability, dielectric loss, capacitance stability, self-healing,
    safety certification, and failure mode are application-specific requirements.
\end{solutionblock}


%%%%%%%%%%%%%%%%%%%%%%%%%%%%%%%%%%%%%%%%%%%%%%%%%%%%%%%%%%%%%%%%%%%%%%%%%%%%%%%
%% Task 5 Condition monitoring and an end-of-life decision
%%%%%%%%%%%%%%%%%%%%%%%%%%%%%%%%%%%%%%%%%%%%%%%%%%%%%%%%%%%%%%%%%%%%%%%%%%%%%%%

\task{Condition monitoring and an end-of-life decision}

A DC-link electrolytic capacitor was characterized when new and again after
several years of operation. The data are listed in
\autoref{table:capacitor_task5_data}.

\begin{table}[ht]
    \centering
    \begin{tabular}{lll}
        \toprule
        Parameter & New condition & Measured condition \\
        \midrule
        Capacitance & \SI{2200}{\micro\farad} & \SI{1760}{\micro\farad} \\
        ESR at the monitoring frequency & \SI{20}{\milli\ohm} & \SI{40}{\milli\ohm} \\
        \bottomrule
    \end{tabular}
    \caption{Condition-monitoring data for Task 5.}
    \label{table:capacitor_task5_data}
\end{table}

The capacitor carries \SI{10}{\ampere} RMS ripple current at an ambient
temperature of \SI{70}{\degreeCelsius}. Its thermal resistance to ambient is
\SI{4}{\kelvin\per\watt}. For this tutorial, use the following maintenance limits:

\begin{itemize}
    \setlength{\itemsep}{-2pt} % Distance between the points (Part1)
    \setlength{\parskip}{-2pt} % Distance between the points (Part2)       
    \item capacitance not below \SI{80}{\percent} of the initial value,
    \item ESR not above twice its initial value,
    \item hot-spot temperature below \SI{95}{\degreeCelsius}.
\end{itemize}

Use $P_\mathrm{ESR}=I_\mathrm{rms}^2R_\mathrm{ESR}, \quad \text{and} \quad
    T_\mathrm{hotspot}=T_\mathrm{amb}+R_\mathrm{th}P_\mathrm{ESR}$.

\subtask{Calculate the percentage capacitance remaining and the ESR increase factor.}

\begin{solutionblock}
    \begin{align}
        \frac{C_\mathrm{meas}}{C_0}
        &=\frac{\SI{1760}{\micro\farad}}{\SI{2200}{\micro\farad}}
        =0.80=\SI{80}{\percent},\\
        \frac{R_\mathrm{ESR,meas}}{R_\mathrm{ESR,0}}
        &=\frac{\SI{40}{\milli\ohm}}{\SI{20}{\milli\ohm}}=2.
    \end{align}
    Both electrical health indicators have reached, but not yet exceeded, the
    defined maintenance limits.
\end{solutionblock}

\subtask{Calculate the capacitor loss and estimated hot-spot temperature in the
new and measured conditions at \SI{10}{\ampere} RMS.}

\begin{solutionblock}
    New condition:
    \begin{align}
        P_0&=(I_\mathrm{rms,ripple})^2 \cdot R_\mathrm{ESR,0}=(\SI{10}{\ampere})^2\cdot\SI{20}{\milli\ohm}
        =\SI{2.0}{\watt},\\
        T_\mathrm{hotspot,0}&=T_\mathrm{ambient} + R_\mathrm{th} \cdot P_0 =\SI{70}{\degreeCelsius}
        +\SI{4}{\kelvin\per\watt}\cdot\SI{2.0}{\watt}
        =\SI{78}{\degreeCelsius}.
    \end{align}
    Measured condition:
    \begin{align}
        P_1&=(I_\mathrm{rms,ripple})^2 \cdot R_\mathrm{ESR,meas}=(\SI{10}{\ampere})^2\cdot\SI{40}{\milli\ohm}
        =\SI{4.0}{\watt},\\
        T_\mathrm{hotspot,1}&=T_\mathrm{ambient} + R_\mathrm{th} \cdot P_1 = \SI{70}{\degreeCelsius}
        +\SI{4}{\kelvin\per\watt}\cdot\SI{4.0}{\watt}
        =\SI{86}{\degreeCelsius}.
    \end{align}
    ESR ageing doubles the internal heating under the same ripple current.
\end{solutionblock}

\subtask{Repeat the measured-condition thermal calculation for a temporary
\SI{13}{\ampere} RMS operating mode.}

\begin{solutionblock}
    \begin{align}
        P_{13\mathrm{A}}&=(I_\mathrm{rms(13),ripple})^2 \cdot R_\mathrm{ESR,meas}=(\SI{13}{\ampere})^2\cdot\SI{40}{\milli\ohm}
        =\SI{6.76}{\watt},\\
        T_\mathrm{hotspot,13\mathrm{A}}&=\SI{70}{\degreeCelsius}
        +\SI{4}{\kelvin\per\watt}\cdot\SI{6.76}{\watt}
        \approx\SI{97.0}{\degreeCelsius}.
    \end{align}
    The high-load mode violates the specified hot-spot limit.
\end{solutionblock}

\subtask{For the same capacitor-current waveform, quantify the change in the
ideal capacitive voltage ripple caused by the capacitance reduction.}

\begin{solutionblock}
    Since ideal capacitive ripple is inversely proportional to capacitance,
    \begin{equation}
        \frac{\Delta V_\mathrm{meas}}{\Delta V_0}
        =\frac{C_0}{C_\mathrm{meas}}
        =\frac{2200}{1760}=1.25.
    \end{equation}
    The ideal capacitive ripple is therefore \SI{25}{\percent} higher than in the
    new condition.
\end{solutionblock}

\subtask{Make a maintenance decision and identify the quantities that should be trended.}

\begin{solutionblock}
    The capacitor can still remain below the thermal limit at the nominal
    \SI{10}{\ampere} operating point, but both capacitance and ESR have reached
    the defined electrical maintenance limits, and the high-load mode is no longer
    thermally acceptable. Planned replacement is therefore justified rather than
    continued unrestricted operation. Capacitance, ESR or dissipation factor,
    leakage current, case/hot-spot temperature, and ripple-current exposure should
    be trended over time.
\end{solutionblock}
